\subsection{An explicit finite algebraic test for elliptic torsion existence}

\begin{proposition}[A bounded test; the output at the first wild gap is not yet evaluated]
\label{prop:B-RD24-finite-test}
\status{Proved; complete proof and original-source audit independently passed by the parent, with an independent rerun of the exact polynomial certificate.  No numerical output for $(p,e,n)=(3,216,3)$ is asserted.}
For every odd prime $p$ and positive integers $e,n$, the existence of a complete discretely valued field of absolute ramification $e$, perfect residue characteristic $p$, and an additive elliptic curve with full rational $p^n$-torsion is equivalent to the finite polynomial test specified below.  The test uses only arithmetic in finite-length Witt vectors, polynomial division and discriminants over $\mathbf Q$, and effective elimination over $\overline{\mathbf F}_p$.  If the test is affirmative, an example exists over a finite extension of $\mathbf Q_p$.
\end{proposition}

\paragraph{Constants and the division polynomial.}
Put $m=p^n$, $N=(m^2-1)/2$, and $a=ne$.  Write
\[
 \psi_m(X)=mX^N+\sum_{i=0}^{N-1}d_iX^i
\]
for the usual odd division polynomial (whose recurrences and torsion-root interpretation are recorded in \cite[section~1.1.3, pp.~8--9]{Burhanuddin2007}) of
\[
 y^2+a_1xy+a_3y=x^3+a_2x^2+a_4x+a_6.
\]
The long-equation version used here can be specified without any division in its coefficient ring.  Put
\[
\begin{gathered}
 b_2=a_1^2+4a_2,\quad b_4=a_1a_3+2a_4,\quad
 b_6=a_3^2+4a_6,\\
 b_8=a_1^2a_6+4a_2a_6-a_1a_3a_4+a_2a_3^2-a_4^2,\qquad
 B(X)=4X^3+b_2X^2+2b_4X+b_6.
\end{gathered}
\]
For odd $j$ let $P_j=\psi_j$, and for even $j$ let
$P_j=\psi_j/\psi_2$.  Their initial values are
\[
\begin{aligned}
 P_0&=0,\quad P_1=P_2=1,\\
 P_3&=3X^4+b_2X^3+3b_4X^2+3b_6X+b_8,\\
 P_4&=2X^6+b_2X^5+5b_4X^4+10b_6X^3+10b_8X^2\\
    &\hspace{7mm}+(b_2b_8-b_4b_6)X+b_4b_8-b_6^2.
\end{aligned}
\]
For subsequent indices the recurrences are
\[
 P_{2r}=P_r(P_{r+2}P_{r-1}^2-P_{r-2}P_{r+1}^2)
\]
and
\[
 P_{2r+1}=
 \begin{cases}
 P_{r+2}P_r^3-B^2P_{r-1}P_{r+1}^3,&r\text{ odd},\\
 B^2P_{r+2}P_r^3-P_{r-1}P_{r+1}^3,&r\text{ even}.
 \end{cases}
\]
They follow from the usual recurrences and the identity $\psi_2^2=B(X)$.
Completing the square and then translating $X$ to remove its quadratic coefficient over $\mathbf Q$ identifies these initial values and recurrences with the short-equation division polynomials in the cited source; thus it also gives their stated torsion-root interpretation.  All the displayed operations are integral polynomial operations.  In particular the $d_i$ can be calculated by this finite recursion.  Define
\[
 c_m=\frac{\operatorname{disc}\bigl(\psi_m(X;0,0,0,-1,0)\bigr)}
                  {64^{\,N(N-1)/6}}\ \in\mathbf Q^\times,
 \qquad t_m=v_p(c_m).
\]
This definition is an exact rational computation, not a conjectured closed formula for the discriminant constant.
Let
\[
 D_0=\begin{cases}6e+11,&p=3,\\11,&p\geq5,\end{cases}
 \qquad
 D_F=a(N-1)(N-2)+e\,t_m+\frac{N(N-1)}6D_0,
\]
and take the deliberately generous positive integer
\[
 D=2|D_F|+a(4N+3)+ND_0+1,\qquad
 H=\max\{2D+1,\,3e+7,\,D_0+1\},\qquad q=\lceil H/e\rceil.
\]
The proof below shows that $D_F$ is in fact nonnegative and bounds the valuation of the monic polynomial discriminant used in the test.  Absolute values in the definition of $D$ merely make positivity manifest before this verification.

\paragraph{The finite coefficient ring.}
All quantified residue-field variables range over $k=\overline{\mathbf F}_p$.
Choose $e$ Witt vectors $c_0,\ldots,c_{e-1}\in W_q(k)$ with the zeroth Witt coordinate of $c_0$ nonzero, and form the finite-length ring
\[
 R=W_q(k)[\pi]\big/\left(\pi^e+pC(\pi)\right),
 \qquad C(\pi)=\sum_{i=0}^{e-1}c_i\pi^i.
\]
This ring is free of rank $e$ over $W_q(k)$, so all its operations are explicit polynomial operations in $eq$ residue-field coordinates.  Its maximal ideal is $(\pi)$ and its length is $eq\geq H$.  Membership in $\pi^rR$ is the finite polynomial condition
\[
                  z\in\pi^rR\quad\Longleftrightarrow\quad
                  (\exists w\in R)\ z=\pi^rw.
\]
Thus every condition below is a finite first-order polynomial formula over $k$.

\paragraph{The test.}
Ask whether there exist such $c_i$, coefficients $a_1,a_2,a_3,a_4,a_6\in R$, an integer $1\leq\delta\leq D_0$, and $N$ pairs $(X_i,Y_i)\in R^2$ satisfying the following conditions.

\begin{enumerate}
\item The universal invariants satisfy
\[
 \Delta\in\pi^\delta R\setminus\pi^{\delta+1}R,
 \qquad c_4^3\in\pi^\delta R.
\]
\item There are no $S,R_1,T_1\in R$ for which all five expressions
\[
\begin{aligned}
 U_1&=a_1+2S,\\
 U_2&=a_2-Sa_1+3r-S^2,\\
 U_3&=a_3+ra_1+2t,\\
 U_4&=a_4-Sa_3+2ra_2-(t+rS)a_1+3r^2-2St,\\
 U_6&=a_6+ra_4+r^2a_2+r^3-ta_3-rta_1-t^2,
\end{aligned}
\qquad r=\pi^{-e}R_1,\quad t=\pi^{-e}T_1
\]
satisfy $\pi^{3e}U_j\in\pi^{3e+j}R$, for $j=1,2,3,4,6$.
This notation contains no inversion of $\pi$ in $R$: first expand each $\pi^{3e}U_j$ and substitute $r=\pi^{-e}R_1,t=\pi^{-e}T_1$; every resulting exponent of $\pi$ is nonnegative.
\item Form the two integral polynomials
\[
\begin{aligned}
 F(X)&=X^N+\sum_{i=0}^{N-1}
       (-C(\pi))^n\pi^{a(N-i-1)}d_iX^i,\\
 G(X,Y)&=Y^2+(a_1\pi^aX+a_3\pi^{2a})Y\\
 &\hspace{9mm}-\bigl(\pi^aX^3+a_2\pi^{2a}X^2
                   +a_4\pi^{3a}X+a_6\pi^{4a}\bigr).
\end{aligned}
\]
For every $i$ require
\[
 F(X_i),G(X_i,Y_i)\in\pi^HR,\qquad
 F'(X_i)G_Y(X_i,Y_i)\notin\pi^{D+1}R.
\]
For $i\ne j$ require $X_i-X_j\notin\pi^{D+1}R$.
\end{enumerate}
Here $R_1$ in condition~2 is a variable, not the coefficient ring $R$.
The test returns \emph{yes} exactly when this finite formula is true in an algebraically closed field of characteristic $p$.

\begin{proof}
We give the bounds and the exact approximation argument, since these distinguish the test from an enumeration of unknown local fields or unknown infinite lifts.
Lemma~\ref{lem:B-RD19} proves that full rational odd-prime torsion excludes additive potentially multiplicative reduction.  Consequently every additive curve in the statement is potentially good, so the following integral-$j$ bounds cover the whole stated problem.

\emph{Minimal discriminants.}
Let an integral minimal equation have potentially good reduction, and put $\delta=v(\Delta)$.  Since $j=c_4^3/\Delta$ is integral, $v(c_4)\geq\delta/3$.  The identity $c_4^3-c_6^2=1728\Delta$ then gives $v(c_6)\geq\delta/2$.  Its short equation has coefficients $-c_4/48$ and $-c_6/864$.  If $p=3$ and $\delta\geq6e+12$, their valuations are at least $4$ and $6$; scaling $x=\pi^2x',y=\pi^3y'$ gives an integral equation with discriminant valuation $\delta-12$, a contradiction.  For $p\geq5$ the same argument applies to $\delta\geq12$.  Hence $\delta\leq D_0$.

An integral equation is nonminimal exactly when a change of variables with scaling factor $\pi$ produces another integral equation.  Indeed, a change with factor $\pi^b$, $b\geq1$, also gives one with factor $\pi$ using the same translations; unit scaling does not affect integrality.  The standard change-of-variables formulas are precisely the $U_j$ above.  The first condition gives $s\in\OO$ because $2$ is a unit.  The second gives $v(r)\geq-v(3)\geq-e$; the third then gives $v(t)\geq-e$.  Conversely these are all the possible translations.  Clearing denominators by $\pi^{3e}$ proves that condition~2 is exactly minimality and depends only on the displayed finite precision.

\emph{The two discriminant bounds.}
Over $\mathbf Q[a_1,a_2,a_3,a_4,a_6,\Delta^{-1}]$, the nonzero $m$-torsion modulo the free involution $P\mapsto-P$ is finite \'etale, with coordinate algebra defined by $\psi_m/m$.  Consequently $\operatorname{disc}(\psi_m)$ is a unit after inverting $\Delta$.  The universal Weierstrass discriminant is irreducible over $\mathbf Q$: after the invertible change to a short equation this is the irreducibility of $4A^3+27B^2$.  Its weighted degree is $12$, while the weighted degree of $\operatorname{disc}(\psi_m)$ is $2N(N-1)$.  It follows that
\[
           \operatorname{disc}(\psi_m)=c_m\Delta^{N(N-1)/6};
\]
specializing to $y^2=x^3-x$ proves exactly the formula defining $c_m$.

Over a lift to a complete DVR, $p/\pi^e=-C(\pi)^{-1}$, so $F(X)$ is exactly
\[
                 \frac{\pi^{aN}}{m}\psi_m(\pi^{-a}X).
\]
It is monic integral, and therefore
\[
 v(\operatorname{disc}F)
   =a(N-1)(N-2)+e\,t_m+\frac{N(N-1)}6\delta
   \leq D_F.
\]
The algebra
\[
                 \mathcal A=\OO[X,Y]/(F,G)
\]
is finite free of rank $2N=m^2-1$, with \'etale generic fibre equal to the coordinate algebra of $E[m]\setminus\{0\}$.
For completeness the needed second resultant has an elementary derivation.  Write
\[
 B(x)=4x^3+b_2x^2+2b_4x+b_6.
\]
After completing the square, at a two-torsion abscissa $\alpha$ the division-polynomial recurrences give
\[
 \psi_m(\alpha)=(-1)^{(m-1)/2}
         f'(\alpha)^{(m^2-1)/4},\qquad
 f(x)=x^3+(b_2/4)x^2+(b_4/2)x+b_6/4.
\]
The base cases are $\psi_1=1$ and $\psi_3(\alpha)=-f'(\alpha)^2$.
For $m=2r+1$ one of the two terms in the odd recurrence vanishes at $\alpha$, and the other proves the formula by induction.  Multiplying over the three roots of $f$ gives
\[
                 \operatorname{Res}(\psi_m,B)
                     =(-1)^{(m-1)/2}\Delta^{N/2}.
\]
The discriminant of the quadratic $G$ is $\pi^{4a}B(\pi^{-a}X)$.
The discriminant formula for a tower of finite free algebras, or its direct determinant proof after splitting $F$, now gives
\[
 v(\operatorname{disc}\mathcal A)
 =2v(\operatorname{disc}F)+a(4N-3)+(N/2)\delta<D.
\]
This proves the required uniform bound without using a formula for $c_m$ that has not been independently established.

\emph{Why approximate roots are exact enough.}
If $E[m]$ is rational, all $N$ roots of the monic $F$ and both roots of the corresponding monic quadratics $G$ lie in $\OO$.  At each point the Jacobian determinant of $(F,G)$ is $F'G_Y$.  The product of these determinants over all $2N$ points is, up to sign, the discriminant of $\mathcal A$.  All the factors are integral, so each has valuation less than $D$.  Distinct roots $X_i,X_j$ have $v(X_i-X_j)\leq v(\operatorname{disc}F)/2<D$, by the Vandermonde product.  Thus reduction modulo the stated precision supplies all the conditions.

Conversely, lift approximate pairs to $\OO$.  Their equation errors have valuation at least $H\geq2D+1$, whereas their Jacobian determinants have valuation at most $D$.  The following elementary multivariate Hensel argument applies to each pair.  If a polynomial map $f$ has integral Jacobian $J$ at $z$, determinant valuation $b\leq D$, and error valuation at least $H$, substitute $z'=z+\pi^{H-b}w$ and consider
\[
                \pi^{-H}\operatorname{adj}(J)
                              f(z+\pi^{H-b}w).
\]
This is an integral polynomial system.  Its linear Jacobian modulo $\pi$ is $(\det J/\pi^b)I$, which is invertible; all nonlinear coefficients are divisible by $\pi$, since $H-2b\geq1$.  The reduction is therefore a linear system with a unique solution, and successive coefficient lifting gives an exact zero.  Since $J$ is invertible over the fraction field, this is also a zero of the original system.  Its difference from $z$ has valuation at least $H-b\geq D+1$.  The distinctness conditions consequently retain $N$ distinct exact abscissas.  Each exact pair is a nonzero $m$-torsion point; its inverse is also rational.  These $2N$ points exhaust $E[m]\setminus\{0\}$.

The first two conditions give a minimal integral equation with positive discriminant valuation and integral $j$.  It is therefore additive and potentially good, and the torsion conclusion follows.

\emph{Why the fields and the decision procedure are bounded.}
By the Cohen structure theorem \cite[Theorem~10.160.8]{Stacks032A}, a mixed-characteristic complete DVR with perfect residue field contains its Witt coefficient ring.  It is generated over that ring by a uniformizer: reduction modulo $p$ and completeness give the basis $1,\pi,\ldots,\pi^{e-1}$, and the relation is Eisenstein of degree $e$.  Conversely every Eisenstein relation used in the test lifts to such a DVR, and $R$ is its quotient modulo $p^q$.  All conditions above are stable under arbitrary lifts because $eq\geq H$.

For the necessity direction, start with its arbitrary perfect residue field $k_0$ and pass to the completed maximal unramified extension.  The absolute ramification remains $e$, rational torsion remains rational, and the new residue field is an algebraic closure $\overline{k_0}$.  The identification of the absolute Galois group of this completed extension with the original inertia group is stated explicitly in \cite[section~5.1.5, pp.~156--157]{Fontaine1994Semi}.  Serre--Tate's good-reduction criterion \cite[Theorem~1, p.~493]{SerreTate1968Good} therefore shows that a bad curve remains bad.  Its integral $j$ remains integral, so it remains additive and potentially good.  Choose a new minimal integral equation after this extension; no invariance claim about the old equation's minimality is needed.

The field $\overline{k_0}$ need not be algebraic over $\mathbf F_p$.  The passage from that field to the field $k=\overline{\mathbf F}_p$ in the test is instead a logical transfer of a finite formula over $\mathbf F_p$.  Addition and multiplication on $W_q$ are universal polynomial operations in $q$ residue-field coordinates, and reduction by the monic Eisenstein polynomial gives polynomial operations on $R$.  Thus the entire displayed test, including its existential and universal quantifiers, is a finite first-order polynomial formula with coefficients in $\mathbf F_p$.  Elimination in algebraically closed fields makes its truth independent of the algebraically closed field of characteristic $p$ used to evaluate it.  In particular a witness over $\overline{k_0}$ gives an affirmative test over $\overline{\mathbf F}_p$, without identifying these two residue fields.

One can describe that elimination procedure directly.  Replace inequations by inverses, express Boolean formulas as constructible sets, and compute existential images by elimination of polynomial ideals.  On each irreducible component of the closure of an image, a Gr\"obner basis over its function field supplies a nonzero denominator whose complement lies in the image; repeat on the remaining closed subset.  Dimension strictly decreases, so this is finite.  Complements handle universal quantifiers.  All coefficients lie in $\mathbf F_p$, so a nonempty constructible output has a point over $\overline{\mathbf F}_p$ with finitely many coordinates in one finite field.  Lifting their Witt coordinates to that finite unramified extension of $\mathbf Q_p$ and the Eisenstein coefficients to finite Witt polynomials produces the asserted finite local-field example.
\end{proof}

\paragraph{Exact computational certificate.}
The program \path{code/verify_finite_elliptic_test.py} reconstructs the division polynomials and computes their exact discriminants and two-torsion resultants for nine specializations, including the degree-$364$ polynomial for $m=27$ on $y^2=x^3-x$.  It also checks $2000$ integer-bound instances.  The original run took $74.857496652$ seconds, and the parent's independent rerun took $74.239452977$ seconds; both passed.  The retained output records the large precision bound at $(3,216,3)$ and explicitly records that its elimination test was not evaluated.  These finite checks support the identities and arithmetic used in the proof; they do not substitute for the universal argument. The formerly open index $216$ is instead decided by the independent construction in Theorem~\ref{thm:B-RD26-wild}.

\paragraph{Scope.}
This bounded test itself has not been evaluated at $(e,n)=(216,3)$. The simplified arithmetic characterization is now proved separately in Theorem~\ref{thm:B-RD26-wild} and Corollary~\ref{cor:B-RD27-exact-three}, without executing this elimination test.  It also does not decide Part~B.4: arbitrary truncated BT data with a finite wild descent action are not thereby supplied with an infinite equivariant BT lift.
